\begin{frame}{Beneficios en renta de personas jurídicas por actividad (2021)}
\begin{table}[h!]
    \centering
    \small
    \setlength{\tabcolsep}{4pt}
    \renewcommand{\arraystretch}{1.05}
    \begin{tabular}{@{}c p{0.62\textwidth} c@{}}
        \toprule
        & \textbf{Actividad económica} & \shortstack{\textbf{Gasto tributario}\\(\% del PIB)} \\
        \midrule
        1  & Transmisión de energía eléctrica & 0,13\% \\
        2  & Actividades de planes de seguridad social obligatoria & 0,11\% \\
        3  & Otras actividades del mercado de valores & 0,11\% \\
        4  & Banco central y bancos comerciales & 0,09\% \\
        5  & Distribución de energía eléctrica & 0,07\% \\
        6  & Extracción de petróleo crudo & 0,06\% \\
        7  & Actividades de administración empresarial & 0,06\% \\
        8  & Seguros de vida & 0,05\% \\
        9  & Producción de malta y cervezas & 0,03\% \\
        10 & Otras actividades de servicio financiero & 0,03\% \\
        \midrule
        & \textbf{Total} & \textbf{0,74\%} \\
        \bottomrule
    \end{tabular}
    \vspace{0.12cm}
    \par\scriptsize{\textit{Fuente: Cuadro 6.3 de El laberinto fiscal de Colombia, con base en Pardo et al. (2024).}}
\end{table}

\note{[Sources] Cuadro 6.3 de El laberinto fiscal de Colombia: https://www.laberintofiscal.co/capitulos/6-inevitables-como-la-muerte/\#table-6-3. Fuente original: Pardo, O.; Heredia, L.; Nieto, A.; y Millán, G. (2024), ``La tributación en impuesto sobre la renta por actividad económica en Colombia'', Coyuntura Económica, vol. 54. El total de 0,74\% del PIB corresponde a la suma de las diez actividades y coincide con la edición definitiva y la versión web.}
\end{frame}
